\chapter{Falsification criteria and boundary cases}

\section{Results that would weaken or reject the central mechanism}

The geometric mechanism would be weakened if:
\begin{enumerate}[leftmargin=1.6em]
\item directly estimated \(\Sop\) channels do not diverge more between contrasting habitats than within habitats;
\item channel divergence is not repeatably associated with divergent selection;
\item curvature-weighted hybrid covariance fails to predict recombinant hybrid fitness out of sample;
\item raw neutral divergence performs as well as or better than the full geometric model;
\item hybrid load persists identically across environments despite the model predicting a large extrinsic component;
\item estimated covariance--curvature alignment has the wrong sign but hybrids nevertheless show the predicted positive load without another mechanism;
\item inferred channel angle is driven entirely by an unstable rank-one choice and disappears for the supported channel dimension.
\end{enumerate}

\section{Biologically meaningful cases in which the mechanism predicts little isolation}

\subsection{Commuting divergence}
\(\G\) and \(\Hh\) change in magnitude but retain common axes. Eigenvalues change; channels do not rotate.

\subsection{Large gap and weak cross-forcing}
The channel is strongly protected and the perturbation barely couples it to other modes.

\subsection{Rotated channels without covariance--curvature coadaptation}
Geometric discrepancy exists, but the trace term is zero or negative.

\subsection{Hybrid variance remains on benign directions}
Segregation does not populate directions strongly penalised by either habitat.

\subsection{Dominance masks early additive load}
F1 hybrids are vigorous even though F2 recombination later reveals a cost.

\section{Scope of the local theory}

The perturbation expansion assumes smooth operator paths and a nonzero gap at the chosen boundary. Near an exact crossing, the correct object is the combined degenerate subspace; rank-one derivatives diverge because the rank-one direction is not identifiable. At large divergence, global geometry, multiple optima, nonquadratic selection, developmental constraints, genotype--environment interactions, and epistasis can all matter.

\begin{summarybox}
The local theory is strongest as an early-speciation model. That is a feature rather than a weakness: it targets the interval in which ecological barriers can arise before many fixed intrinsic incompatibilities exist.
\end{summarybox}
